\chapter{Placenta problems}
\index{Placenta problems}

\medskip
\noindent\textit{Educational information; seek urgent maternity advice for bleeding or reduced fetal movement.}

\section*{What the placenta does}
The placenta connects the pregnant person and fetus, allowing oxygen and nutrients to pass and waste products to leave. Problems can involve placental position, separation, abnormal attachment, blood flow, or function.

\section*{Important conditions}
Placenta praevia is a low placenta that covers or approaches the cervix and commonly causes painless bleeding later in pregnancy. Placental abruption is early separation from the uterus and may cause bleeding, abdominal or back pain, contractions, or fetal distress. Placenta accreta spectrum is abnormal attachment that can cause severe bleeding during birth. Reduced placental function may affect fetal growth.

\section*{Assessment and management}
Assessment may include maternal and fetal examination, ultrasound, blood tests, fetal heart monitoring, and repeat growth or blood-flow scans. Management depends on gestational age, bleeding, fetal wellbeing, placental position, and the specific problem. Care may involve monitoring, hospital admission, medicines, specialist planning, or delivery.

\section*{When to seek emergency care}
Contact the maternity unit or local emergency service immediately for any vaginal bleeding, severe or persistent abdominal or back pain, contractions, dizziness or fainting, fluid loss, reduced fetal movement, or feeling very unwell. Do not wait for a routine appointment.

\section*{Sources}
\begin{itemize}
\item ACOG bleeding during pregnancy: \url{https://www.acog.org/womens-health/faqs/bleeding-during-pregnancy}
\item ACOG placenta accreta spectrum: \url{https://www.acog.org/clinical/clinical-guidance/obstetric-care-consensus/articles/2018/12/placenta-accreta-spectrum}
\end{itemize}
